% Manually maintained resume for AI-company applications.
\documentclass[11pt,letterpaper]{article}

\usepackage[T1]{fontenc}
\usepackage{lmodern}
\usepackage[top=0.50in,bottom=0.50in,left=0.70in,right=0.70in]{geometry}
\usepackage{enumitem}
\usepackage{microtype}
\usepackage{xcolor}
\usepackage[hidelinks]{hyperref}
\hypersetup{pdftitle={Junzhao Yang - Resume},pdfauthor={Junzhao Yang}}

\definecolor{accent}{HTML}{26736B}
\definecolor{muted}{HTML}{4F5963}
\definecolor{rulecolor}{HTML}{D8DEE2}

\pagestyle{empty}
\setlength{\parindent}{0pt}
\setlength{\parskip}{0pt}
\renewcommand{\familydefault}{\sfdefault}

\newcommand{\ressection}[1]{%
  \par\vspace{6pt}
  {\color{accent}\bfseries\fontsize{10.5}{12.5}\selectfont\MakeUppercase{#1}}\par
  \vspace{2pt}{\color{rulecolor}\hrule height 0.5pt}\vspace{4pt}
}

\newcommand{\resdate}[1]{{\small\color{muted}#1}}

\newcommand{\publication}[3]{%
  \item \begin{tabular*}{\linewidth}[t]{@{}p{\dimexpr\linewidth-2.1in\relax}@{\extracolsep{\fill}}p{2in}@{}}
    \raggedright\textbf{#1} & \raggedleft\color{accent}\bfseries #2
  \end{tabular*}\par
  {\footnotesize\color{muted}#3}
}

\begin{document}

{\fontsize{20}{23}\selectfont\bfseries Junzhao Yang}\par
\vspace{3pt}
{
\small Pittsburgh, PA \enspace\textbar\enspace
+1 (412) 844-1653
  \enspace\textbar\enspace \href{mailto:junzhaoy@andrew.cmu.edu}{junzhaoy@andrew.cmu.edu}
  \enspace\textbar\enspace \href{https://fizzydavid.github.io/}{fizzydavid.github.io}}\par

\ressection{Education}
\textbf{Carnegie Mellon University}, Computer Science Department
\hfill \resdate{2024--Present}\par
Ph.D. in Computer Science (expected 2027). Advisor: Richard Peng \par
Research areas: numerical linear algebra, spectral graph theory \par
\vspace{3pt}
\textbf{Tsinghua University}, Institute for Interdisciplinary Information Sciences
\hfill \resdate{2020--2024}\par
B.S. in Computer Science, \href{https://iiis.tsinghua.edu.cn/en/yaoclass/}{Yao Class}.\par

\ressection{Projects and Experience}
\textbf{FizzyProver} \hfill \resdate{May 2026--Present}\par
{\small\color{muted}Sole developer \textperiodcentered\ Multi-agent system for mathematical research}\par
\begin{itemize}[leftmargin=1.25em,topsep=2pt,itemsep=1pt,parsep=0pt]
  \small
  \item Designed a recursive Codex workflow that schedules independent proof subproblems in parallel and alternates prover and verifier passes to critique and revise proof attempts.
  \item Built and deployed a web platform used by 20+ academic users across 300+ runs. Among the runs we have examined to date, including runs powered by GPT-5.5, we have not observed any incorrect final proofs; by contrast, our direct GPT-5.5 Pro attempts produced flawed proofs.
\end{itemize}

\textbf{Megvii Internship} \hfill \resdate{Winter 2021}\par
{\small\color{muted}Research Intern}\par
\begin{itemize}[leftmargin=1.25em,topsep=2pt,itemsep=1pt,parsep=0pt]
  \small
  \item Reproduced a graph-neural-network post-training algorithm from a research paper in PyTorch.
  \item Optimized its large-matrix computations for GPU memory efficiency; portions of the implementation were later reused by the team.
\end{itemize}

\ressection{Honors and Awards}
\begin{itemize}[leftmargin=1.25em,topsep=1pt,itemsep=1pt,parsep=0pt]
  \item ICPC North America Championship 2025, 2nd place, Silver Medal.
  \item \textbf{46th ICPC World Finals, 10th place, Bronze Medal} (held in 2024).
  \item \textbf{International Olympiad in Informatics 2019, 5th place, Gold Medal}.
\end{itemize}
\vspace{0.3pt}
\begin{itemize}[leftmargin=1.25em,topsep=0pt,itemsep=0pt,parsep=0pt]
  \item Yao Award, Bronze Medal (awarded to five Yao Class students in 2023).
\end{itemize}

\ressection{Selected Publications}
{\footnotesize\color{muted}Authors are listed alphabetically, as is customary in theoretical computer science.}\par
\begin{enumerate}[label={[\arabic*]},leftmargin=2.1em,topsep=3pt,itemsep=1pt,parsep=0pt]
  \publication{Fast Convergence of Gradient Descent on Random $\ell_p$-Norm Regression}{FOCS 2026 (to appear)}{Yang Liu, Richard Peng, Weina Wang, \textbf{Junzhao Yang}}
  \publication{Approximate Spanning Tree Counting from Uncorrelated Edge Sets}{FOCS 2026 (to appear)}{Yang P. Liu, Richard Peng, \textbf{Junzhao Yang}}
  \publication{Entrywise Approximate Solutions for SDDM Systems in Almost-Linear Time}{STOC 2026}{Angelo Farfan, Mehrdad Ghadiri, \textbf{Junzhao Yang}}
  \publication{Entrywise Approximation for Matrix Inversion and Linear Systems}{SODA 2026}{Mehrdad Ghadiri, Hoai-An Nguyen, \textbf{Junzhao Yang}}
  \publication{Numerical Linear Algebra in Linear Space}{SODA 2026}{Yiping Liu, Hoai-An Nguyen, \textbf{Junzhao Yang}}
  \publication{The Cost of Parallelizing Boosting}{SODA 2024}{Xin Lyu, Hongxun Wu, \textbf{Junzhao Yang}}
  \publication{Tight Time-Space Lower Bounds for Constant-Pass Learning}{FOCS 2023}{Xin Lyu, Avishay Tal, Hongxun Wu, \textbf{Junzhao Yang}}
\end{enumerate}

\end{document}
